%
% LaTeX-Rahmen f�r Arbeiten am Lehrstuhl Hegering
%
% Harald Roelle, 2001, 2002
%
% basierend auf Arbeiten von Helmut Reiser, Boris Gruschke und Stephen Heilbronner
%


%Link im PDF
\ifpdf
  \pdfbookmark[0]{Titel}{Titel}%
\fi

%%%%%%%%%%%%%%%%%%%%%%%%%%%%%%%
% erste Seite

\thispagestyle{empty}

\begin{center}

\vspace*{-2cm}

{\Huge INSTITUT F�R INFORMATIK\\[1mm]}
DER LUDWIG--MAXIMILIANS--UNIVERSIT�T M�NCHEN\\

\vspace*{1cm}

\includegraphics[width=0.3\textwidth]{lmu_siegel}

\vspace*{2cm}

{\Large \textbf{\typeOfThesis}}\\

\vspace{2.0cm}
{\Huge \textbf{\titleOfThesisOne}}\\
\vspace*{3mm}
{\Huge \textbf{\titleOfThesisTwo}}\\
\vspace*{3mm}
{\Huge \textbf{\titleOfThesisThree}}\\

\vspace{1.5cm}

{\LARGE \authorOfThesis}

%\vspace{2cm}

%\parbox{1cm}{
%  \begin{large}
%    \begin{tabbing}
%      Aufgabensteller:
%        \=\aufgabensteller\\[2mm]
%      Betreuer: 
%        \>\betreuerOne\\
%        \>\betreuerTwo\\
%        \>\betreuerThree\\[5mm]
%      Abgabetermin: 
%        \> \abgabeTagZahl.~\abgabeMonatText~\abgabeJahrZahl\\
%    \end{tabbing}
%  \end{large}
%}\\

\end{center}

\newpage

%%%%%%%%%%%%%%%%%%%%%%%%%%%%%%%
% zweite Seite

\thispagestyle{empty}
\cleardoublepage 

%%%%%%%%%%%%%%%%%%%%%%%%%%%%%%%
% dritte Seite (Kopie der ersten)

\thispagestyle{empty}

\begin{center}

\vspace*{-2cm}

{\Huge INSTITUT F�R INFORMATIK\\[1mm]}
DER LUDWIG--MAXIMILIANS--UNIVERSIT�T M�NCHEN\\

\vspace*{1cm}

\includegraphics[width=0.3\textwidth]{lmu_siegel}

\vspace*{2cm}

{\Large \textbf{\typeOfThesis}}\\

\vspace{2.0cm}
{\Huge \textbf{\titleOfThesisOne}}\\
\vspace*{3mm}
{\Huge \textbf{\titleOfThesisTwo}}\\
\vspace*{3mm}
{\Huge \textbf{\titleOfThesisThree}}\\

\vspace{1.5cm}

{\LARGE \authorOfThesis}

\vspace{2cm}

\parbox{1cm}{
  \begin{large}
    \begin{tabbing}
      Aufgabensteller:
        \=\aufgabensteller\\[2mm]
      Betreuer: 
        \>\betreuerOne\\
        \>\betreuerTwo\\
        \>\betreuerThree\\[5mm]
      Abgabetermin: 
        \> \abgabeTagZahl.~\abgabeMonatText~\abgabeJahrZahl\\
    \end{tabbing}
  \end{large}
}\\

\end{center}
